\documentclass[11pt,reqno]{amsart}
\usepackage[T1]{fontenc}
\usepackage[utf8]{inputenc}
\usepackage{amsmath,amssymb,amsthm,mathtools}
\usepackage[margin=1.1in]{geometry}
\usepackage{booktabs}
\usepackage[protrusion=true,expansion=false]{microtype}
\usepackage[colorlinks=true,linkcolor=blue,citecolor=blue,urlcolor=blue]{hyperref}

\theoremstyle{plain}
\newtheorem{theorem}{Theorem}[section]
\newtheorem{lemma}[theorem]{Lemma}
\newtheorem{proposition}[theorem]{Proposition}
\newtheorem{corollary}[theorem]{Corollary}
\theoremstyle{definition}
\newtheorem{definition}[theorem]{Definition}
\newtheorem{remark}[theorem]{Remark}
\newtheorem{example}[theorem]{Example}
\newtheorem{problem}[theorem]{Problem}

\newcommand{\Z}{\mathbb{Z}}
\newcommand{\Q}{\mathbb{Q}}
\newcommand{\R}{\mathbb{R}}
\newcommand{\OK}{\mathcal{O}_K}
\newcommand{\OO}{\mathcal{O}}
\newcommand{\OF}{\mathcal{O}_F}
\newcommand{\rr}{\mathfrak{r}}
\newcommand{\pp}{\mathfrak{p}}
\newcommand{\Gcl}{\mathcal{G}}
\newcommand{\SoS}{\Sigma\OK^{2}}
\newcommand{\tp}{\succ 0}
\DeclareMathOperator{\Var}{Var}
\DeclareMathOperator{\Tr}{Tr}
\DeclareMathOperator{\Nm}{N}
\DeclareMathOperator{\disc}{disc}
\DeclareMathOperator{\dist}{dist}
\newcommand{\Fm}{F_m}

\begin{document}

\title[The class $\Gcl$ beyond degree five]
{The class $\Gcl$ beyond degree five: exact witnesses, the $2$-unramified classification
in degrees $6$--$8$, and the bottom of the variance spectrum}
\author{}
\date{\today}

\begin{abstract}
Let $\Gcl$ be the class of totally real fields $K$ with $2\OK^{\tp}\subseteq\Sigma\OK^{2}$. Its
members of degree at most $5$ are known (eight fields), and it is expected that there are no others.
This note records a computational and theoretical attack on degrees $6$, $7$, $8$ and on the
structure of small-variance fields.

(i) We give an exact decision procedure for ``$2\alpha\in\Sigma\OK^{2}$'' and a search for
\emph{exceptional} elements (witnesses), and run it on every totally real field of degree $3$--$8$
in the LMFDB lists ($193{,}467$ cubic, $13{,}073$ quartic, $9{,}494$ quintic, $9{,}815$ sextic,
$301$ septic and $920$ octic fields). In degrees $\le5$ it reproduces the Kala--Yatsyna
classification exactly; in degrees $6$, $7$ and $8$ every field has an explicit, independently
certified witness.

(ii) We show that for fields in which $2$ is unramified, membership in $\Gcl$ implies every
necessary condition (budget--cost, sharpened rigidity, square dichotomy, and their relative
versions) used in the classification of totally real fields of degree $6$--$8$ with a universal
$\Z$-form. Combining this with that classification and with explicit $\Gcl$-witnesses for all the
fields that it has to construct ($18$ sextic and $394$ octic fields; it constructs no septic
field), we obtain:
\emph{no totally real field of degree $6$, $7$ or $8$ in which $2$ is unramified lies in $\Gcl$.}

(iii) We determine the bottom of the variance spectrum. All totally real algebraic integers with
$\Var(x)\le2$ are listed in degrees $\le10$, and an auxiliary-function (linear programming)
certificate, verified in interval arithmetic, shows that every other totally real algebraic integer
has $\Var(x)\ge 1.69$. Consequently the values of $\Var$ below $1.69$ are exactly
$0,\tfrac54,\tfrac{14}9,\tfrac{27}{16}$, and a totally real field has $\nu(K)<1.69$ if and only if it
contains $\Q(\sqrt5)$ or $\Q(\zeta_7)^{+}$. As a by-product,
$\liminf_{\deg x\to\infty}\Var(x)\ge1.69$, which improves the bound $\sqrt e=1.6487\ldots$ quoted
in the working notes.

(iv) We record shape and arithmetic-statistics data from the census, explicit witnesses for the
seven smallest non-member real cyclotomic fields left open by the dyadic criterion ($m=36,44,52,76,100,108,116$),
and a reduction of the remaining problem (sextic fields with $2$ ramified) to explicitly described
cases.
\end{abstract}

\maketitle
\tableofcontents

\section{Introduction}\label{sec:intro}

For a totally real field $K$ of degree $d$ write $\OK^{\tp}$ for its totally positive integers and
$\SoS$ for the sums of squares. Put
\[
 \Gcl=\{K\ \text{totally real}:\ 2\OK^{\tp}\subseteq\SoS\}.
\]
Kala and Yatsyna \cite{KY} proved that the members of degree at most $5$ are
\begin{equation}\label{eq:G5}
 \Q,\ \Q(\sqrt2),\ \Q(\sqrt3),\ \Q(\sqrt5),\ \Q(\sqrt2,\sqrt5),\ \Q(\zeta_{20})^{+},\ \Q(\zeta_7)^{+},\
 \Q(\rho)\quad(\rho^3=4\rho+2),
\end{equation}
and it is expected that \eqref{eq:G5} is all of $\Gcl$. The working notes \cite{PB,PD,PR} develop
the \emph{trace variance} $\Var(x)=\tau(x^{2})-\tau(x)^{2}$, $\tau=\frac1d\Tr_{K/\Q}$, as the natural
height for this problem, and the companion repository \cite{TwoO} pushes the related \emph{lifting
problem} (universal $\Z$-forms) through degrees $6$, $7$ and $8$ using the same invariants. The
purpose of this note is to transfer as much as possible of that machinery to $\Gcl$, to test it
against complete field tables, and to understand the fields of small variance, which are the only
places where members of $\Gcl$ can live.

Throughout, $\Lambda_K=\OK/\Z$ with the positive definite form $\Var$ is the \emph{variance lattice}
(up to scaling it is the shape of $K$), and
\[
 \nu(K)=\min_{x\in\OK\setminus\Z}\Var(x),\qquad
 \nu_2(K)=\min_{x\in(\Z+\rr)\setminus\Z}\Var(x),\qquad
 t(K)=\min_{\alpha\in\OK^{\tp}\setminus\Z}\tau(\alpha),
\]
where $\rr=\prod_{\pp\mid2}\pp^{\lceil e_\pp/2\rceil}$ is the dyadic square-root radical; one has
$2\OK\subseteq\rr$, $\nu\le\nu_2\le4\nu$, and $\nu_2=4\nu$ when $2$ is unramified.

\subsection*{Results}
\begin{enumerate}
\item \emph{Anatomy of a representation} (\S\ref{sec:anatomy}). A complete description of what a
representation $2\alpha=\sum x_i^2$ looks like in the low-trace range, valid without any
hypothesis on $2$. It contains the budget--cost inequality and rigidity of \cite{PB} and adds a
\emph{half-square dichotomy}: if $\tau(\alpha)<\nu(K)$ and $K\in\Gcl$ then $2\alpha-c$ is a square
in $\OK$ for an integer $0\le c<2\min_j\sigma_j(\alpha)$.
\item \emph{An exact witness search} (\S\ref{sec:algo}) and a census of all LMFDB fields of degree
$3$--$8$ (\S\ref{sec:census}); Theorem~\ref{thm:census}.
\item \emph{The $2$-unramified classification in degrees $6$--$8$} (\S\ref{sec:unram});
Theorem~\ref{thm:unram}.
\item \emph{The bottom of the variance spectrum} (\S\ref{sec:spectrum}); Theorems~\ref{thm:lp}
and~\ref{thm:smallnu}.
\item \emph{Shapes and statistics} (\S\ref{sec:stats}), \emph{residual cyclotomic fields}
(\S\ref{sec:cyclo}), the \emph{ramified sextic problem} (\S\ref{sec:ramified}), and corrections to
the working notes (\S\ref{sec:corrections}).
\end{enumerate}

All computations use PARI/GP~2.15.4 \cite{PARI}; the code and all outputs are in the repository
accompanying this note (\S\ref{sec:repro}). Every exclusion of a field rests on an explicit element
whose exceptional status is certified by exact arithmetic, independently of the (GRH-conditional)
class group and unit computations used to \emph{find} it.

\section{Anatomy of a representation}\label{sec:anatomy}

\begin{lemma}\label{lem:anatomy}
Let $K$ be totally real, $\alpha\in\OK^{\tp}\setminus\Z$, and $2\alpha=\sum_{i=1}^{N}x_i^{2}$ with
$x_i\in\OK\setminus\{0\}$. Let $S=\{i:x_i\notin\Z\}$. Then $S\neq\emptyset$, $\sum_ix_i\in\rr$, and
\begin{equation}\label{eq:budget}
 2\tau(\alpha)=\sum_{i\in S}\bigl(\Var(x_i)+\tau(x_i)^{2}\bigr)+\sum_{i\notin S}x_i^{2}.
\end{equation}
Moreover:
\begin{enumerate}
\item If $|S|=1$, say $S=\{1\}$, put $c=\sum_{i\ge2}x_i^{2}\in\Z_{\ge0}$. Then $2\alpha-c=x_1^{2}$,
$c<2\min_j\sigma_j(\alpha)$, $x_1\in(c+\rr)\setminus\Z$, and $2\tau(\alpha)\ge\nu_2(K)+c$.
\item If $|S|=2$, say $S=\{1,2\}$, and $\tau(\alpha)<\frac12\nu_2(K)$, then $x_2=\varepsilon x_1+k$
with $\varepsilon=\pm1$, $k\in\Z$, and
\[
 4\alpha=(2x_1+\varepsilon k)^{2}+\kappa,\qquad \kappa=k^{2}+2\sum_{i\ge3}x_i^{2}\in\Z_{\ge0},\quad
 \kappa\equiv0,3\pmod 4,
\]
with $\tau(\alpha)\ge\Var(x_1)+\kappa/4$. If $\kappa=0$ then $\alpha=x_1^{2}$.
\item If $|S|\ge3$ then $\tau(\alpha)\ge\frac32\nu(K)$.
\end{enumerate}
\end{lemma}

\begin{proof}
If $S=\emptyset$ then $2\alpha\in\Z$, so $\alpha\in\Z$. The coupling $\sum_ix_i\in\rr$ is
\cite[dyadic coupling lemma]{PB}: $s=\sum x_i$ has $s^{2}\equiv\sum x_i^{2}\equiv0\pmod{2\OK}$, so
$v_\pp(s)\ge\lceil e_\pp/2\rceil$ for every $\pp\mid2$. Identity \eqref{eq:budget} is $\tau$ of the
representation, with $\tau(x^{2})=\Var(x)+\tau(x)^{2}$.

(1) $x_1^{2}=2\alpha-c$ is totally positive, so $c<2\sigma_j(\alpha)$ for all $j$. For integers
$n\equiv n^{2}\pmod 2$, hence $\sum_{i\ge2}x_i\equiv c\pmod{2}$; as $2\in\rr$, the coupling gives
$x_1\equiv-\sum_{i\ge2}x_i\equiv c\pmod\rr$. Thus $x_1\in(\Z+\rr)\setminus\Z$, so
$\Var(x_1)\ge\nu_2$, and \eqref{eq:budget} gives $2\tau(\alpha)\ge\Var(x_1)+c$.

(2) Put $c_0=\sum_{i\ge3}x_i\in\Z$. The coupling gives $x_1+x_2\in-c_0+\rr$, and $2x_2\in\rr$ gives
$x_1-x_2\in-c_0+\rr$. If neither $x_1+x_2$ nor $x_1-x_2$ is rational, both have variance
$\ge\nu_2$, and the parallelogram law for the quadratic form $\Var$ gives
$2\nu_2\le\Var(x_1+x_2)+\Var(x_1-x_2)=2\Var(x_1)+2\Var(x_2)\le2\tau(x_1^2)+2\tau(x_2^2)\le4\tau(\alpha)$,
contrary to the hypothesis. Hence $x_2=\varepsilon x_1+k$. Expanding,
$2\alpha=2x_1^{2}+2\varepsilon kx_1+k^{2}+M$ with $M=\sum_{i\ge3}x_i^{2}$, so
$4\alpha=(2x_1+\varepsilon k)^{2}+k^{2}+2M$. Since $k^{2}+M=2(\alpha-x_1^{2}-\varepsilon kx_1)\in2\OK\cap\Z$,
$M\equiv k^{2}\pmod2$ and $\kappa\equiv3k^{2}\pmod4$. Finally
$\tau(\alpha)=\tau\bigl((x_1+\varepsilon k/2)^{2}\bigr)+\kappa/4\ge\Var(x_1)+\kappa/4$, and $\kappa=0$
forces $k=0=M$, i.e.\ $\alpha=x_1^{2}$.

(3) Three summands of \eqref{eq:budget} are each $\ge\nu$.
\end{proof}

\begin{corollary}[$\Gcl$-dichotomies]\label{cor:dich}
Let $K\in\Gcl$, $K\ne\Q$, and $\alpha\in\OK^{\tp}\setminus\Z$. Write $\nu=\nu(K)$, $\nu_2=\nu_2(K)$.
\begin{enumerate}
\item \textup{(budget--cost \cite[budget--cost inequality]{PB})} $\tau(\alpha)\ge\min(\nu,\frac12\nu_2)$.
\item \textup{(half-square dichotomy)} If $\tau(\alpha)<\nu$, then $2\alpha-c=y^{2}$ for some $c\in\Z$
with $0\le c<2\min_j\sigma_j(\alpha)$ and some $y\in(c+\rr)\setminus\Z$; moreover
$2\tau(\alpha)\ge\nu_2+c$.
\item \textup{(square window)} If $\tau(\alpha)<\min(\frac32\nu,\ \nu+\frac34,\ \frac12\nu_2)$, then
$\alpha$ is a square in $\OK$.
\item \textup{(witness by trace)} If $\tau(\alpha)<\min(\nu,\frac12\nu_2)$ for some
$\alpha\in\OK^{\tp}\setminus\Z$ of an arbitrary totally real $K$, then $2\alpha\notin\SoS$; in
particular $K\notin\Gcl$, and $\alpha$ itself is a witness.
\end{enumerate}
\end{corollary}

\begin{proof}
(1), (2), (4) are immediate from Lemma~\ref{lem:anatomy}. For (3): $|S|\ge3$ is excluded by
$\tau<\frac32\nu$, and $|S|=1$ by $\tau<\frac12\nu_2$ (Lemma~\ref{lem:anatomy}(1)); so $|S|=2$, and
$\nu+\kappa/4\le\tau(\alpha)<\nu+\frac34$ forces $\kappa\in\{0,1,2\}$, hence $\kappa=0$ by the
congruence, hence $\alpha=x_1^{2}$.
\end{proof}

\begin{remark}\label{rem:rung}
Apply (2) to the \emph{lower rung} $\alpha_-=x-n$ of any $x\in\OK\setminus\Z$, where
$n=\lceil\theta_1\rceil-1$ and $\theta_1$ is the least conjugate of $x$. Then
$\min_j\sigma_j(\alpha_-)=\theta_1-n\in(0,1]$, so $c\in\{0,1\}$: \emph{if $\tau(x-n)<\nu$ then
$2x-2n-c$ is a square in $\OK$ for some $c\in\{0,1\}$}. This is the $\Gcl$-analogue of the ``rung
split'' of \cite[degree-$8$ rung split]{TwoO}, and it is the only way in which the lower rung can be cheap. The same
holds for the upper rung $\lfloor\theta_d\rfloor+1-x$.
\end{remark}

\subsection{The unramified case, and relative versions}
If $2$ is unramified then $\rr=2\OK$, $(\Z+\rr)/\Z=2\Lambda_K$, and $\nu_2=4\nu$. Then
Corollary~\ref{cor:dich}(1) reads $\tau(\alpha)\ge\nu$, the half-square alternative never occurs
(it forces $\tau(\alpha)\ge2\nu$), and the window of (3) is $\min(\frac32\nu,\nu+\frac34)$.

Let $F\subsetneq K$ be a subfield, $\Var_{K/F}$ the relative variance (an $F$-valued totally
nonnegative quadratic form vanishing exactly on $F$), and
$\nu_{K/F}=\min\{\tau_F(\Var_{K/F}x):x\in\OK\setminus\OF\}$; by the law of total variance
\cite[ANOVA identity]{PB}, $\tau(x^{2})\ge\tau_F(\Var_{K/F}x)$.

\begin{proposition}[relative version, $2$ unramified]\label{prop:rel}
Let $K\in\Gcl$ with $2$ unramified in $K$, $F\subsetneq K$, and $\alpha\in\OK^{\tp}\setminus F$.
Then $\tau(\alpha)\ge\nu_{K/F}$. If moreover $\tau(\alpha)<\frac32\nu_{K/F}$, then
$\alpha=y^{2}+\gamma y+\beta$ with $y\in\OK\setminus\OF$, $\beta,\gamma\in\OF$,
$\tau_F(\Var_{K/F}y)\le\tau(\alpha)$, and $4\beta-\gamma^{2}$ totally nonnegative.
\end{proposition}

\begin{proof}
Let $S_F=\{i:x_i\notin F\}\ne\emptyset$. If $|S_F|\ge2$ then $\tau(\alpha)\ge\nu_{K/F}$; if
$|S_F|\ge3$ then $\tau(\alpha)\ge\frac32\nu_{K/F}$. If $S_F=\{1\}$, the coupling (with $\rr=2\OK$)
gives $x_1\in\OF+2\OK$, so $x_1=f+2w$ with $w\notin F$ and $\tau(x_1^{2})\ge4\nu_{K/F}$, i.e.\
$\tau(\alpha)\ge2\nu_{K/F}$. If $S_F=\{1,2\}$, then $x_1\pm x_2\in\OF+2\OK$, and if both are
outside $F$ the parallelogram law for $\tau_F\circ\Var_{K/F}$ gives $\tau(\alpha)\ge2\nu_{K/F}$.
Otherwise $x_2=\pm x_1+\gamma$ with $\gamma\in\OF$; then
$\alpha=x_1^{2}\pm\gamma x_1+\beta$ with $2\beta=\gamma^{2}+\sum_{i\ge3}x_i^{2}$, $\beta\in\OK\cap F=\OF$,
and $4\beta-\gamma^{2}=\gamma^{2}+2\sum_{i\ge3}x_i^{2}\succeq0$.
\end{proof}

\section{An exact decision procedure, and certificates}\label{sec:algo}

\subsection{Deciding $\beta\in\SoS$}
Let $\beta\in\OK^{\tp}$. Every summand $x$ of a representation $\beta=\sum x_i^{2}$ satisfies
$\sigma_j(x)^{2}\le\sigma_j(\beta)$ for all $j$, hence lies in the finite set
\[
 X(\beta)=\Bigl\{x\in\OK\setminus\{0\}:\ \textstyle\sum_j\sigma_j(x)^{2}/\sigma_j(\beta)\le d\Bigr\}
 \cap\{x:\beta-x^{2}\succeq0\}.
\]
We enumerate the ellipsoid by Fincke--Pohst (PARI's \texttt{qfminim} on the weighted Gram matrix,
LLL-reduced, at a working precision adapted to the spread of the conjugates of $\beta$) and test
$\beta-x^{2}\succeq0$ exactly. A depth-first search over multisets, with a memo on remainders and a
direct look-up for one and two squares, decides whether $\beta\in\SoS$. Every representation found
is re-verified by exact multiplication.

\subsection{Rigorous signs}
All decisions about signs of conjugates are exact. The embedding matrix is computed to $\ge100$
digits, and a numerical sign is accepted only when the computed value exceeds an explicit bound
$(\|x\|_1+1)\max|M_{ij}|\,10^{-90}$ for the error; otherwise the sign is decided from a certified
high-precision evaluation (precision doubled until $|v|$ exceeds a rigorous evaluation error) or by a
Sturm sequence on the characteristic polynomial. This matters: in large-regulator fields the
totally positive generators of small ideals are unavoidably unbalanced (conjugates of size
$10^{\pm50}$ occur for cubic fields of discriminant $\approx2\cdot10^{6}$), and a first version of
the code that used a fixed tolerance produced wrong sign vectors of fundamental units there. The
bug was caught by the final certification step below.

\subsection{The witness search}
For a field $K$ we try, in order: (U) the \emph{unit criterion}: every $u\in U_K^{+}$ with
$u,2u\notin K^{\times2}$ is a witness, since $\Nm(u)=1$ allows at most two squares and equality in
O'Meara's superadditivity forces $u$ or $2u$ to be a square \cite[unit dichotomy lemma]{PB}; (T) all
$\alpha\in\OK^{\tp}\setminus\Z$ with $\tau(\alpha)\le T$, enumerated in $\Lambda_K$ using
$\Var(\alpha)\le(d-1)\tau(\alpha)^{2}$ and the least totally positive shift of each class
(Corollary~\ref{cor:dich}(4) explains why such elements are natural witnesses); (N) for every ideal
of norm $\le N_{\max}$ that is narrowly principal, every totally positive generator modulo $U_K^{2}$,
reduced in trace by squares of units.

\subsection{Certificates}
A witness $\alpha$ is accepted only after (a) total positivity is re-proved by a Sturm sequence of
its characteristic polynomial, with no floating point, and (b) the decision procedure is rerun from
scratch on $2\alpha$ and finds no representation. Neither step uses the class group or the units, so
every exclusion below is unconditional, although the search that found the witness used
\texttt{bnfinit} (correct under GRH).

\section{The census in degrees $3$--$8$}\label{sec:census}

We ran the search on the LMFDB lists of totally real fields \cite{LMFDB} as mirrored in
\cite{Visser} (the lists used by the \texttt{sails} project). The bounds were $T=3$ ($T=2$ in degree
$3$) and $N_{\max}=64,128,300,300,400$ in degrees $4,5,6,7,8$ ($N_{\max}=64$ in degree $3$); the few
survivors were rerun with $N_{\max}=1500$.

\begin{theorem}\label{thm:census}
Every one of the 9815 fields of degree 6 ($d_K\le1.00\cdot10^{8}$), 301 fields of degree 7 ($d_K\le2.14\cdot10^{8}$), 920 fields of degree 8 ($d_K\le6.97\cdot10^{9}$) in the LMFDB lists has an explicit element $\alpha\in\OK^{\tp}$ with $2\alpha\notin\SoS$, certified as in \S\ref{sec:algo}; none of these fields lies in $\Gcl$. In degrees 3, 4, 5 the fields without a witness are exactly the members of \eqref{eq:G5} in the lists.

\end{theorem}

\begin{table}[t]
\caption{The census. U, T, N: number of fields whose first witness is a unit, a small-trace element,
or a small-norm generator. ``max $N$'' is the largest norm of a first witness of type N.}
\label{tab:census}
\begin{tabular}{rrrrrrrrr}
\toprule
$d$ & fields & $\max d_K$ & U & T & N & left & $h^+=h$ & max $N$\\
\midrule
3 & 192468 & $3.37\cdot10^{6}$ & 111535 & 1 & 80930 & 2 (49,148) & 42.1\% & 197\\
4 & 13073 & $1.00\cdot10^{6}$ & 7113 & 1118 & 4840 & 2 (1600,2000) & 30.7\% & 81\\
5 & 9494 & $1.00\cdot10^{7}$ & 5170 & 3097 & 1227 & 0 & 45.5\% & 17\\
6 & 9815 & $1.00\cdot10^{8}$ & 5024 & 4413 & 378 & 0 & 40.6\% & 43\\
7 & 301 & $2.14\cdot10^{8}$ & 156 & 145 & 0 & 0 & 48.2\% & --\\
8 & 920 & $6.97\cdot10^{9}$ & 405 & 515 & 0 & 0 & 51.1\% & --\\
\bottomrule
\end{tabular}

\end{table}

\begin{remark}[validation against \cite{KY}]
In degrees $3$, $4$, $5$ the procedure finds a witness for every field except exactly the members
in \eqref{eq:G5} (discriminants $49,148$; $1600,2000$; none), which survive up to $N_{\max}=1500$.
Since our implementation shares no code with \cite{KY}, this is an independent check of their
classification on these lists, and conversely their theorem is a check of our code: a single false
witness or a missed representation in a member would have shown up here.
\end{remark}

\begin{remark}[the extreme cases]
The fields that need the largest witnesses are imprimitive fields with a subfield of very small
variance and a large discriminant: the seven quartic fields with $d_K\in[5.7\cdot10^{5},10^{6}]$
that survive $N_{\max}=64$ all contain $\Q(\sqrt5)$ or $\Q(\sqrt2)$, have normalised shape minimum
$\approx0.09$ (against a median $0.58$), and need witnesses of norm $71$--$81$. This is the
``inherited witness'' phenomenon of \cite[``What remains, and why'']{PB} seen in data: a small-variance subfield blocks
every trace-based witness, and only norm-based witnesses remain, whose size grows with $d_K$.
\end{remark}

\section{Fields with $2$ unramified in degrees $6$, $7$, $8$}\label{sec:unram}

The repository \cite{TwoO} proves that no totally real field of degree $6$, $7$ or $8$ admits a
universal $\Z$-form. Its proof uses the following necessary conditions on a field $K$ (we refer to
the notes in \cite{TwoO} for the precise statements):
\begin{itemize}
\item[(BC)] $\tau(\alpha)\ge\nu(K)$ for all $\alpha\in\OK^{\tp}\setminus\Z$, and its integer-part
consequences (downspread, the quadratic filter);
\item[(RG)] rigidity: $\tau(\alpha)<\frac32\nu(K)$ implies $\alpha=y^{2}+cy+b$ with $y\notin\Z$,
$\Var(y)\le\tau(\alpha)$, $\kappa=4b-c^{2}\ge0$, $\kappa\equiv0,3\pmod4$; in particular the square
dichotomy $\tau(\alpha)<\min(\frac32\nu,\nu+\frac34)\Rightarrow\alpha\in\OK^{2}$;
\item[(REL)] the relative versions of (BC) and (RG) over every subfield;
\end{itemize}
and, only at the very end, condition (N) at explicit elements.

\begin{proposition}\label{prop:transfer}
If $K\in\Gcl$ and $2$ is unramified in $K$, then $K$ satisfies (BC), (RG) and (REL).
\end{proposition}

\begin{proof}
(BC) and (RG) are Corollary~\ref{cor:dich} with $\nu_2=4\nu$ together with
Lemma~\ref{lem:anatomy}(2) (which gives exactly the shape $y^{2}+cy+b$ with
$\kappa=k^{2}+2\sum_{i\ge3}x_i^{2}\ge0$); (REL) is Proposition~\ref{prop:rel}.
\end{proof}

\begin{theorem}\label{thm:unram}
No totally real field of degree $6$, $7$ or $8$ in which $2$ is unramified belongs to $\Gcl$.
\end{theorem}

\begin{proof}
By Proposition~\ref{prop:transfer}, such a field passes every test of \cite{TwoO} that precedes
the final condition-(N) step. That computation leaves, in degree $6$, the $18$ fields of
\cite[degree-$6$ Case B]{TwoO} (Case B; Case A is empty); in degree $7$, nothing; in degree $8$, the four
fields of Case A and the $390$ fields of Case B over $\Q(\sqrt5)$. For each of these
$18+4+390$ fields we found an explicit $\alpha\in\OK^{\tp}$ with $2\alpha\notin\SoS$, certified as
in \S\ref{sec:algo} (files \texttt{results/g2o\_d6.txt}, \texttt{g2o\_d8A.txt},
\texttt{g2o\_d8B.txt}); Table~\ref{tab:unram} lists the ones with $2$ unramified in degree~$6$ and
the two octic fields that survive all tests of \cite{TwoO} except (N).
\end{proof}

\begin{table}[t]
\caption{$\Gcl$-witnesses for the fields left by the pipeline of \cite{TwoO} that have $2$
unramified, in degree $6$, and for the two octic fields excluded there only by (N).}
\label{tab:unram}
\begin{tabular}{rll}
\toprule
$d_K$ & witness $\alpha$ (minimal polynomial) & $\Nm_{K/\Q}\alpha$\\
\midrule
300125 & $x^6 - 14x^5 + 73x^4 - 182x^3 + 227x^2 - 133x + 29$ & 29\\
485125 & $x^6 - 13x^5 + 64x^4 - 151x^3 + 177x^2 - 96x + 19$ & 19\\
820125 & $x^3 - 6x^2 + 9x - 3$ & 9\\
1416125 & $x^6 - 14x^5 + 75x^4 - 191x^3 + 232x^2 - 117x + 19$ & 19\\
3162625 & $x^6 - 15x^5 + 86x^4 - 234x^3 + 299x^2 - 142x + 4$ & 4\\
\midrule
442050625 & $x^4 - 9x^3 + 27x^2 - 31x + 11$ & 121\\
656505625 & $x^4 - 9x^3 + 27x^2 - 31x + 11$ & 121\\
\bottomrule
\end{tabular}

\end{table}

\begin{remark}
(a) The proof uses the computations of \cite{TwoO} as a black box; we have checked that each
exclusion test there is one of (BC), (RG), (REL), and reproduced its degree-$6$ and degree-$7$
runs (\S\ref{sec:repro}). (b) Most of the fields in the list have $2$ ramified; they are not
covered by Proposition~\ref{prop:transfer}, but our witnesses exclude them from $\Gcl$ anyway.
(c) For $\Q(\sqrt5,\zeta_7^{+})$ ($d_K=300125$) and the two octic fields, the $\Gcl$-witness is the
same element at which (N) fails in \cite{TwoO}: the element of trace $14$ and norm $29$, and the
element of trace $18$ in $K$ (trace $9$ and norm $11$ in the quartic subfield of discriminant $725$),
respectively.
\end{remark}

\section{The bottom of the variance spectrum}\label{sec:spectrum}

\subsection{Enumeration}
Up to $x\mapsto\pm x+k$ and conjugation, a totally real algebraic integer of degree $e$ is a
real-rooted monic $f\in\Z[t]$ with $a_1\in\{0,-1,\dots,-\lfloor e/2\rfloor\}$, and
$e\Var(x)=\frac{e-1}{e}a_1^{2}-2a_2$. We enumerated all such $f$ with $\Var\le2$ (Robinson's method,
the enumerator of \cite{TwoO} in its unfiltered mode) for $e\le10$ and kept the irreducible ones.

\begin{table}[t]
\caption{The variance spectrum below $1.9$ in degrees $\le10$: value, degrees in which it occurs,
number of polynomials (up to $x\mapsto\pm x+k$).}
\label{tab:spectrum}
\begin{tabular}{llll}
\toprule
$\Var$ & decimal & degrees & \#\\
\midrule
$5/4$ & 1.2500 & 2 & 1\\
$14/9$ & 1.5556 & 3 & 1\\
$27/16$ & 1.6875 & 4 & 1\\
$7/4$ & 1.7500 & 4,8 & 3\\
$44/25$ & 1.7600 & 5 & 1\\
$16/9$ & 1.7778 & 9 & 1\\
$9/5$ & 1.8000 & 10 & 3\\
$65/36$ & 1.8056 & 6 & 2\\
$90/49$ & 1.8367 & 7 & 3\\
$46/25$ & 1.8400 & 5,10 & 9\\
$37/20$ & 1.8500 & 10 & 12\\
$119/64$ & 1.8594 & 8 & 8\\
$152/81$ & 1.8765 & 9 & 15\\
$17/9$ & 1.8889 & 6 & 4\\
$189/100$ & 1.8900 & 10 & 43\\
\bottomrule
\end{tabular}

\end{table}

\subsection{Auxiliary functions}
Let $I=[s_0,s_1]\subseteq[0,\frac12]$, $Q_1,\dots,Q_k\in\Z[t]$ and $c_1,\dots,c_k\ge0$ be such that
\begin{equation}\label{eq:aux}
 f_I(t)=\dist(t,I)^{2}-\sum_ic_i\log|Q_i(t)|\ \ge\ m_I\qquad(t\in\R).
\end{equation}
If $x$ is a totally real algebraic integer with $\tau(x)\in I$ and $Q_i(x)\ne0$ for all $i$, then
averaging \eqref{eq:aux} over the conjugates of $x$ and using $|\Nm(Q_i(x))|\ge1$ gives
$\Var(x)\ge\frac1e\sum_j\dist(\theta_j,I)^{2}\ge m_I$. Every $x$ can be normalised by $x\mapsto\pm x+k$
to have $\tau(x)\in[0,\frac12]$.

We took for the $Q_i$ the linear polynomials $t-k$ and all integer shifts and reflections, with
roots in $[-4.5,5]$, of the irreducible polynomials of degree $\le7$ found in the enumeration
($959$ polynomials), solved the linear program for the $c_i$ on a grid, rounded the $c_i$ down to
rationals, lowered each $m_I$ by $5\cdot10^{-4}$, and \emph{verified} \eqref{eq:aux} on all of
$\R$, for each of the $10$ intervals $I=[\frac{k}{20},\frac{k+1}{20}]$ covering $[0,\frac12]$: on $[-7,7.5]$ by interval arithmetic
(\texttt{mpmath.iv}, $120$ bits) with adaptive subdivision, using
$\log|Q_i(t)|\le\log\max_{\text{cell}}|Q_i|$; outside, by the explicit increasing minorant
$(|t|-\frac12)^{2}-\sum_ic_i(\log\|Q_i\|_1+\deg Q_i\log|t|)$.

\begin{theorem}\label{thm:lp}
Let $m^{*}=\frac{33839}{20000}=1.69195$. There is an explicit set $\mathcal Q$ of $109$ irreducible
polynomials in $\Z[t]$, of degrees $\le7$ (file \texttt{results/lp/cert\_deg7.json}), such that every
totally real algebraic integer $x$ which is not, after $x\mapsto\pm x+k$ $(k\in\Z)$, a root of a
member of $\mathcal Q$ satisfies
\[
 \Var(x)\ \ge\ m^{*}.
\]
The members of $\mathcal Q$ have degree $\le7$, so their roots appear in the enumeration, and among
them only the following have variance $<m^{*}$: $\frac{1+\sqrt5}2$ ($\Var=\frac54$),
$2\cos\frac{2\pi}7$ ($\frac{14}9$), and a root of $t^{4}-t^{3}-3t^{2}+t+1$ ($\frac{27}{16}$; it generates
the totally real quartic field of discriminant $725$, the one of least discriminant). Consequently
\[
 \{\Var(x):x\ \text{a totally real algebraic integer}\}\cap[0,m^{*})=\{0,\tfrac54,\tfrac{14}9,\tfrac{27}{16}\},
\]
each value being attained by a unique integer up to $x\mapsto\pm x+k$ and conjugation, and
$\Var(x)\ge m^{*}$ for every totally real algebraic integer of degree $\ge5$. With the $36$
polynomials of degree $\le5$ of \texttt{cert\_deg5.json} one already gets the bound
$\frac{32597}{20000}=1.62985$.
\end{theorem}

\begin{theorem}\label{thm:smallnu}
Let $K$ be totally real and $m^{*}$ as in Theorem~\ref{thm:lp}. Then $\nu(K)<m^{*}$ if and only if
$K$ contains $\Q(\sqrt5)$ or $\Q(\zeta_7)^{+}$; and then $\nu(K)=\frac54$ if $\sqrt5\in K$, and
$\nu(K)=\frac{14}9$ otherwise. In particular $\nu(K)\ge m^{*}$ for every totally real field without
a quadratic or cubic subfield (for instance every primitive field of degree $\ge4$, and every field
of degree prime to $6$).
\end{theorem}

\begin{proof}
Let $x\in\OK\setminus\Z$ realise $\nu(K)$. If $\nu(K)<m^{*}$, Theorem~\ref{thm:lp} identifies $x$, up
to $x\mapsto\pm x+k$ and conjugation, as the golden ratio, $2\cos\frac{2\pi}7$, or a root of
$t^{4}-t^{3}-3t^{2}+t+1$. The last generates the quartic field $F_{725}$ of discriminant
$725=5^{2}\cdot29$, which contains $\Q(\sqrt5)$; so in that case already $\nu(K)\le\frac54$, and the
value $\frac{27}{16}$ never occurs as $\nu(K)$. Conversely $\Q(\sqrt5)$ and $\Q(\zeta_7)^{+}$ contain
elements of variance $\frac54$ and $\frac{14}9$, and $\frac54$ is the absolute minimum
\cite[sharp floor]{PB}.
\end{proof}

\begin{remark}
The exhaustive enumeration also lists the fields generated by integers of variance $\le2$ (file
\texttt{results/spectrum/fields\_Vle2.txt}): there are $2,2,4,4,18,43,195,463,734$ of them in degrees
$2,\dots,10$ (of which $2,2,1,4,7,43,150,453,650$ primitive). Any totally real field with
$\nu(K)\le2$ whose minimal vector has degree $e\le10$ contains one of them. Their number grows
quickly with the degree, in line with $\liminf\Var\le\lambda_{\mathrm{SSS}}<2$.
\end{remark}

\begin{remark}
(a) Theorem~\ref{thm:lp} gives $\liminf_{\deg x\to\infty}\Var(x)\ge m^{*}$; the working notes
\cite[the limiting-variance theorem]{PB} quote $\sqrt e=1.6487\ldots$ (Schur's argument via Hermite discriminants;
see also \cite{Dubickas,CY}), and the upper bound $\lambda_{\mathrm{SSS}}\le1.8984$ comes from
Smith's realisation theorem \cite{Smith,OSS}. We have not found a better lower bound in the
literature, but have not made an exhaustive search. (b) The degree-$\le5$ certificate already gives
$m^{*}=1.6298$, enough to show that $\frac54$ and $\frac{14}9$ are the only nonzero values below
$1.6298$; this sharpens the ``bottom of the spectrum'' theorem of \cite[bottom of the spectrum]{PB}, which
proves that $(\frac54,\frac{14}9)$ contains no value.
(c) The quartic field of discriminant $725$ is precisely the subfield on which the two octic
survivors of \cite{TwoO} fail (N) and $\Gcl$ (Table~\ref{tab:unram}): small-variance fields are
where the hard cases live.
\end{remark}

\subsection{Consequences for $\Gcl$}
The eight members \eqref{eq:G5} have $\nu\in\{\frac54,\frac{14}9,2,\frac{20}9,3\}$. By
Theorem~\ref{thm:smallnu}, a field $K\in\Gcl$ with $\nu(K)<m^{*}$ contains $\Q(\sqrt5)$ or
$\Q(\zeta_7)^{+}$; every other member has
$\nu(K)\ge m^{*}$, and by the degree-free exclusion \cite[degree-free exclusion]{PB} its minimal vector has degree
$e$ with $m^{*}\le\nu(K)\le2e+2\sqrt{e^{2}-1}$.

\section{Shapes and arithmetic statistics}\label{sec:stats}

The lattice $(\Lambda_K,\Var)$ is the shape of $K$ up to scaling \cite{MSRG}, with determinant
$d_K/d^{d}$. We record the \emph{normalised shape minimum}
$\hat\nu(K)=\nu(K)\,(d^{d}/d_K)^{1/(d-1)}\in(0,\gamma_{d-1}]$.

\begin{table}[t]\caption{Normalised shape minimum $\hat\nu$ in the census: quantiles, the Hermite bound $\gamma_{d-1}$, medians by type of first witness, and the share of fields with $h^+=h$ (equivalently $U_K^+=U_K^2$).}\label{tab:shape}
\begin{tabular}{rrrrrrrrrrr}\toprule
$d$ & $q_{0.01}$ & $q_{0.1}$ & $q_{0.5}$ & $q_{0.9}$ & max & $\gamma_{d-1}$ & U & T & N & $h^+=h$\\\midrule
4 & 0.08 & 0.12 & 0.58 & 0.90 & 1.19 & 1.26 & 0.55 & 0.58 & 0.62 & 30.7\%\\
5 & 0.47 & 0.55 & 0.73 & 0.94 & 1.20 & 1.41 & 0.73 & 0.72 & 0.80 & 45.5\%\\
6 & 0.27 & 0.34 & 0.74 & 0.95 & 1.32 & 1.52 & 0.74 & 0.77 & 0.37 & 40.6\%\\
7 & 0.85 & 0.88 & 0.97 & 1.08 & 1.25 & 1.67 & 0.96 & 0.98 & -- & 48.2\%\\
\bottomrule\end{tabular}\end{table}

Table~\ref{tab:shape} summarises the shapes. Three features stand out.
(a) In the prime degrees $5$ and $7$ the shape minima concentrate in a band around $0.7$--$1.0$, well
below the Hermite bound, while in degrees $4$, $6$, $8$ there is a heavy lower tail: these are the
imprimitive fields, whose minimal vector lies in a subfield of small variance and whose shape is
therefore skewed. (b) In degree $6$, witnesses of type T (small trace, i.e.\ $t(K)$ small
compared with $\nu$ and $\nu_2$: Corollary~\ref{cor:dich}(4)) occur for well-rounded shapes and
witnesses of type N for skewed ones: the medians of $\hat\nu$ are $0.77$ and $0.37$. This is the
budget--cost mechanism seen statistically: a round shape forces $\nu$ to be large compared with the
cheapest totally positive element, while a skewed (imprimitive) shape inherits a small $\nu$ from a
subfield and only a norm-type witness is left. In degree $4$ the distinction is not visible.
(c) Between one half and two thirds of the fields have $U_K^{+}\neq U_K^{2}$; in odd degree all of
these are excluded by the unit criterion, since then $(2)$ is not the square of an ideal and $2u$
can never be a square. The proportion of fields with $h^{+}=h$ that we observe ($31$--$48\%$ in
degrees $4$--$7$) is the kind of statistic predicted by the heuristics of Dummit and Voight for
signature ranks of units \cite{DV}; we have not compared the numbers with their predictions.


\section{The residual real cyclotomic fields}\label{sec:cyclo}

The dyadic criterion of \cite{PD} leaves, among $m=4m_0$, the family with $-1\in\langle2\rangle$
modulo $m_0$ and $\varphi(m_0)/\operatorname{ord}_{m_0}(2)$ odd; for $m_0\le99$ these are
$m_0=3$, $5$, $9$, $11$, $13$, $19$, $25$, $27$, $29$, $37$, $43$, $53$, $59$, $61$, $67$, $81$, $83$
(we re-derived this table). The first two are
members. Our search gives explicit witnesses for the next ones:

\begin{center}\begin{tabular}{rrrrrl}\toprule $m$ & $m_0$ & $[F_m:\Q]$ & type & $\Nm\alpha$ & $\Tr\alpha$\\\midrule
36 & 9 & 6 & trace & 24 & 18\\
44 & 11 & 10 & norm & 89 & 28\\
\bottomrule\end{tabular}\end{center}


These exclusions are unconditional: each witness is certified as in \S\ref{sec:algo}, with no
class-group or GRH input. They are residual for the dyadic criterion of \cite{PD}; the working notes
\cite{PB} state that the conductors $4q$ with $q\le37$ are settled by the signature law and that all
conductors not of the form $4q$ are settled as well, so for these $m$ our witnesses are an
independent, explicit confirmation rather than new exclusions. ($F_{116}$ has $h^{+}/h=16$, so it is
also excluded by the unit criterion of \cite{PB}, and indeed its witness is a unit.) For $F_{52}$ the witness has norm $13$:
it is a totally positive generator of the ramified prime above $13$, as predicted by the signature
law ($13\equiv1\pmod4$ and $4\mid\operatorname{ord}_{13}2=12$), and the exact decision procedure
replaces the analytic non-representation criterion, which needs $q\ge269$.

\section{The sextic problem with $2$ ramified}\label{sec:ramified}

By Theorem~\ref{thm:unram} and the census, a sextic member of $\Gcl$ would have $2$ ramified and
discriminant beyond the LMFDB list. We record how the minimal-vector method splits in that case.
Let $x_1$ realise $\nu$, with rungs $\alpha_\pm$ as in Remark~\ref{rem:rung}.

\emph{Regime 1: $\nu_2\ge2\nu$.} Then Corollary~\ref{cor:dich}(1) gives $\tau(\alpha)\ge\nu$ for
every $\alpha\in\OK^{\tp}\setminus\Z$, so the downspread bound $d\nu\ge2(\nu-1)^{2}$ and the
quadratic filter hold, and Case A of the enumeration of \cite{TwoO} (with its $15{,}720$
irreducible candidates, $\nu\le4.79$) covers all such $K$ with $\Q(x_1)=K$. What is lost is
rigidity, whose window shrinks to $\tau<\frac12\nu_2$ when $\nu_2<3\nu$; the $9{,}785$ candidates
that \cite{TwoO} excludes by rigidity must instead be tested at the level of the field.

\emph{Regime 2: $\nu\le\nu_2<2\nu$.} Budget--cost only gives $\tau\ge\frac12\nu_2$. A cheap rung is
then a \emph{half-square}: $2x_1-2n-c=y^{2}$ with $c\in\{0,1\}$ (Remark~\ref{rem:rung}). If one rung
is regular, the downspread argument gives $\frac32\nu\le\sqrt{2d\nu}+2$, i.e.\ $\nu\le7.77$; if both
are half-squares, $y^{2}+z^{2}=N$ for an integer $N\in[2\nu,4\nu)$ and $\nu\le15.75$. In both cases
one enumerates $y$ with $p_2(y)=d(2\tau(\alpha_-)-c)$ and $p_4(y)=d\bigl(4\nu+(2\tau(\alpha_-)-c)^{2}\bigr)$
fixed, as in the rung split of \cite[degree-$8$ rung split]{TwoO}.

\emph{Case B} (the minimal vector lies in a proper subfield $F$). If $[K:F]$ is odd, an element of
$F$ that is a square in $K$ is already a square in $F$ (otherwise $F(\sqrt\beta)\subseteq K$ would
have even degree over $F$). Hence every half-square condition of Corollary~\ref{cor:dich}(2) for
elements of $F$ must already hold in $F$, with $\nu=\nu(K)=\nu(F)$.

\begin{proposition}\label{prop:quadB}
Let $K\in\Gcl$, and suppose that a minimal vector of $\Lambda_K$ lies in a quadratic subfield $F$
with $[K:F]$ odd. Then $F\in\{\Q(\sqrt2),\Q(\sqrt3),\Q(\sqrt5)\}$.
\end{proposition}

\begin{proof}
Write $F=\Q(\sqrt D)$, $D>1$ squarefree, so $\nu(K)=\nu(F)$ and every $\alpha\in\OF^{\tp}\setminus\Z$
with $\tau(\alpha)<\nu(F)$ must have $2\alpha-c\in\OF^{2}$ for some $c\in\{0,1,\dots\}$ with
$c<2\min\sigma(\alpha)$. Take $\alpha$ to be the lower rung of the minimal vector, so $c\in\{0,1\}$
(Remark~\ref{rem:rung}).

If $D\equiv1\pmod4$, then $\nu=D/4$, the minimal vector is $\omega=\frac{1+\sqrt D}2$ up to
$\pm\omega+k$, and $\alpha=\omega+u$ with $u=\lceil\frac{\sqrt D-1}2\rceil$ has
$\tau(\alpha)=u+\frac12<\frac{\sqrt D}2+1<\frac D4$ as soon as $D\ge13$ (note $\min\sigma(\alpha)=u+\frac{1-\sqrt D}2\in(0,1]$). But
$2\alpha-c=(2u+1-c)+\sqrt D$ is not a square in $\OF$: if $\bigl(\frac{p+q\sqrt D}2\bigr)^{2}$ has
$\sqrt D$-coefficient $1$ then $pq=2$, and then $p\not\equiv q\pmod 2$, so $\frac{p+q\sqrt D}2\notin\OF$.
Hence $D=5$.

If $D\equiv2,3\pmod4$, then $\nu=D$, the minimal vector is $\sqrt D$, and $\alpha=\sqrt D+u$ with
$u=\lfloor\sqrt D\rfloor+1$ has $\tau(\alpha)=u<D$ for $D\ge3$. A square $(p+q\sqrt D)^{2}$ with
$\sqrt D$-coefficient $2$ has $pq=1$, so $2\alpha-c=2u-c+2\sqrt D$ is a square only if
$2u-c=D+1$, i.e.\ $D=2\lfloor\sqrt D\rfloor+1-c\le2\sqrt D+1$, i.e.\ $D\le5$. This leaves $D=2$
(where $\tau(\alpha)=\nu$) and $D=3$ (where $2\alpha=(1+\sqrt3)^{2}$ with $c=0$).
\end{proof}

The three fields of Proposition~\ref{prop:quadB} are exactly the quadratic members of $\Gcl$. A
direct computation (file \texttt{src/quadB.gp}) confirms that for squarefree $D\le2000$ only
$D=2,3,5$ satisfy the half-square condition for \emph{all} $\alpha\in\OF^{\tp}\setminus\Z$ with
$\tau(\alpha)<\nu(F)$. If $[K:F]=2$ (a cubic $F$ in a sextic $K$), the half-square condition no longer
descends, but it is still strong: if some $\alpha\in\OF^{\tp}$ with $\tau(\alpha)<\nu(F)$ has no
half-square in $F$, then $K=F\bigl(\sqrt{2\alpha-c}\bigr)$ for one of the finitely many admissible
$c$, so $K$ is determined by $F$ up to finitely many choices.

\begin{proposition}\label{prop:cubicB}Let $K\in\Gcl$ be sextic, with a minimal vector of $\Lambda_K$ in a cubic subfield $F$. Then $\nu(F)=\nu(K)\le9.58$ (both rungs cost at least $\frac12\nu$, and the spread $s$ of the minimal vector satisfies $s^{2}\le6\nu$), hence $d_F\le(3\nu)^{3}\le23{,}700$, and $t(F)\ge\frac12\nu(F)$. Among the 214 cubic fields with $\nu(F)\le9.58$ (all have $d_F\le11777$), 143 violate $t(F)\ge\frac12\nu(F)$; for 68 of them some $\alpha\in\OF^{\tp}$ with $\tau(\alpha)<\nu(F)$ has no half-square in $F$, which forces $K=F(\sqrt{2\alpha-c})$ for an admissible $c$; the 92 fields so obtained all have certified witnesses. The remaining cubic fields, for which $K$ is not determined, are exactly $d_F=49$, $d_F=81$, $d_F=148$ (that is, $\Q(\zeta_7)^{+}$, $\Q(\zeta_9)^{+}$ and $\Q(\rho)$).\end{proposition}

\begin{proposition}\label{prop:reg1}No sextic $K\in\Gcl$ is in regime~$1$ ($\nu_2\ge2\nu$) with a minimal vector generating $K$.\end{proposition}

\begin{proof}By the discussion above $K$ is generated by a root of one of the $9{,}785$ polynomials that \cite{TwoO} excludes only by rigidity. They define 6429 distinct fields (2856 with $2$ ramified), and each has a certified witness (file \texttt{results/regime1.txt}).\end{proof}

\begin{proposition}\label{prop:reg2X}No sextic $K\in\Gcl$ has a minimal vector $x_1$ generating $K$ with both rungs regular.\end{proposition}\begin{proof}If both rungs are regular, the downspread inequalities hold and $\nu\le4.79$, so the enumeration of \cite{TwoO} without its quadratic filter (which needs $\tau\ge\nu$ for \emph{all} totally positive $\alpha$, not available in regime~$2$) contains the minimal polynomial of $x_1$. This enumeration (\texttt{src/enum6g.c}, \texttt{src/run\_regime2X.sh}) has $178{,}438$ leaves, exactly as in \cite{TwoO}, and outputs $42{,}232$ polynomials. Of the 16666 irreducible ones, 126 are excluded by BC0, 66 are excluded by HS, 1294 are excluded by NS, 5826 are excluded by R1, 9273 are excluded by UNR, 81 are excluded by W, where UNR means $2$ unramified (Theorem~\ref{thm:unram}), R1 means $\nu_2\ge2\nu$ (Proposition~\ref{prop:reg1}), NS that $x_1$ is not minimal in $\OK$, BC0 an element with $\tau<\min(\nu,\frac12\nu_2)$ (a witness by Corollary~\ref{cor:dich}(4)), HS an element with $\tau<\nu$ and no half-square (Corollary~\ref{cor:dich}(2)), and W a certified witness.\end{proof}

\begin{proposition}\label{prop:reg2Y}No sextic $K\in\Gcl$ has a minimal vector $x_1$ generating $K$ with a half-square rung.\end{proposition}\begin{proof}After $x_1\mapsto\pm x_1+k$ we may assume that the lower rung is a half-square and $\tau(x_1)\in[-\frac12,\frac12]$. Then $y^{2}=2x_1-C$ ($C=2n+c$, $c\in\{0,1\}$); if $\Tr x_1=-a_1$ and $\sum\theta_j^{2}=a_1^{2}-2a_2$, the power sums $p_2(y)=-2a_1-6C$ and $p_4(y)=4(a_1^{2}-2a_2)+4Ca_1+6C^{2}$ are fixed, the least $y_j^{2}$ lies in $(0,2-c]$, $\frac12(\nu+c)\le\tau(x_1-n)<\nu$, $\Var(y)\ge\nu$, both rungs of $y$ and the upper rung of $x_1$ cost at least $\frac12\nu$, and $\nu\le15.75$; if $\nu>7.77$ the upper rung is a half-square too, so $|y_j|<2\sqrt\nu$. The enumeration of all such $y$ (mode \texttt{GYL} of \texttt{src/enum6g.c}, validated against brute force in degree $4$ for $\nu\le8$ and in degree $6$ for $\nu\le\frac83$, see \texttt{src/validate/}) visits 422{,}110{,}544 leaves and returns 40{,}374 polynomials $x_1$. Of the 27907 irreducible ones, 13050 are excluded by BC0, 8174 are excluded by HS, 6667 are excluded by NS, 16 are excluded by W (notation as in Proposition~\ref{prop:reg2X}).\end{proof}


\subsection{Summary for degree six}
\begin{theorem}\label{thm:sextic}
Let $K\in\Gcl$ be a totally real sextic field, and let $x_1$ be a minimal vector of $\Lambda_K$.
Then $2$ is ramified in $K$, $K$ is not among the $9{,}815$ sextic fields of the LMFDB list, and one
of the following holds.
\begin{enumerate}
\item \textup{(Case A, regime 2)} $\Q(x_1)=K$, $\nu_2(K)<2\nu(K)$, and at least one rung of $x_1$ is a
half-square ($2x_1-2n-c=y^{2}$ with $y\in\OK$, $c\in\{0,1\}$); $\nu(K)\le7.77$ if the other rung is
regular and $\nu(K)\le15.75$ if both are half-squares.
\item \textup{(Case B, quadratic)} $F=\Q(x_1)\in\{\Q(\sqrt2),\Q(\sqrt3),\Q(\sqrt5)\}$ and
$\nu(K)=\nu(F)\in\{2,3,\frac54\}$.
\item \textup{(Case B, cubic)} $F=\Q(x_1)\in\{\Q(\zeta_7)^{+},\Q(\zeta_9)^{+},\Q(\rho)\}$, $[K:F]=2$, and
$\nu(K)=\nu(F)\in\{\frac{14}9,2,\frac{20}9\}$.
\end{enumerate}
\end{theorem}

\begin{proof}
$2$ is ramified by Theorem~\ref{thm:unram}, and $K$ is not in the list by Theorem~\ref{thm:census}.
If $\Q(x_1)=K$ and $\nu_2\ge2\nu$, Proposition~\ref{prop:reg1} applies; if both rungs are
regular, Proposition~\ref{prop:reg2X} applies; otherwise Remark~\ref{rem:rung} and the downspread
bounds above give (1). If $\Q(x_1)$ is quadratic,
$[K:\Q(x_1)]=3$ is odd and Proposition~\ref{prop:quadB} gives (2). If it is cubic,
Proposition~\ref{prop:cubicB} gives (3).
\end{proof}

The members of $\Gcl$ appearing in (2) and (3) are $\Q(\sqrt2),\Q(\sqrt3),\Q(\sqrt5),\Q(\zeta_7)^{+},
\Q(\rho)$; only $\Q(\zeta_9)^{+}$ is not a member. So, apart from regime~$2$ of Case~A, a sextic member
of $\Gcl$ would have to be an extension of a member of $\Gcl$ of degree $2$ or $3$ which inherits
its variance minimum, with $2$ ramified: precisely the ``inherited'' situation that the census
identifies as extreme (Table~\ref{tab:inherited}).

\begin{problem}\label{prob:sextic}
Close the two remaining cases of Theorem~\ref{thm:sextic}.
\begin{enumerate}
\item[(1)] \emph{Half-square rungs.} Enumerate $y$ with $y^{2}=2x_1-C$, $C=2n+c$, $c\in\{0,1\}$, as in the
rung split of \cite{TwoO}: if $x_1$ has $\Tr x_1=-a_1$ and $\sum\theta_j^{2}=a_1^{2}-2a_2$, then
\[
 p_2(y)=-2a_1-6C,\qquad p_4(y)=4(a_1^{2}-2a_2)+4Ca_1+6C^{2}
\]
are fixed, $\frac12(\nu_2+c)\le\tau(x_1-n)<\nu$, the least $y_j^{2}$ lies in $(0,2-c]$, and $y$ itself
obeys the downspread with the weaker cost $\frac12\nu$. The modification of the enumerator is
straightforward; the search is larger than in \cite{TwoO} because $p_2(y)$ reaches about $190$ when both
rungs are half-squares.
\item[(2)] \emph{Inherited minima.} Bound the relative variance of $K$ over $F\in\{\Q(\sqrt2),\Q(\sqrt3),
\Q(\sqrt5)\}$ ($[K:F]=3$) and over $F\in\{\Q(\zeta_7)^{+},\Q(\zeta_9)^{+},\Q(\rho)\}$ ($[K:F]=2$) when $2$
ramifies in $K$. The relative budget--cost of Proposition~\ref{prop:rel} then has a relative half-square
alternative ($2\alpha-\gamma\in\OK^{2}$ with $\gamma\in\Sigma\OF^{2}$), which costs a factor $2$ in the
relative downspread bound; the relative enumerations of \cite{TwoO} have to be redone with that bound.
\end{enumerate}
Together with Theorem~\ref{thm:unram} this would complete the classification of sextic members of
$\Gcl$.
\end{problem}

\section{Remarks on the working notes}\label{sec:corrections}
We checked several statements of \cite{PB,PD} against independent computations.
\begin{enumerate}
\item The values $\nu,\nu_2,t$ for the members, for $\Q(\sqrt{13})$, and for the cyclic cubic fields
of conductor $13,91,117,133$ in \cite[shape section]{PB} agree with ours.
\item The table of excluded and surviving $m_0\le99$ in \cite[table of $m_0\le99$]{PD} agrees with ours.
\item The effective Northcott subsection of \cite{PB} states that the cubic scan with $d_K\le3000$ leaves
$\{49,81,169,229,257,316,321,404,469,473,564\}$ ``together with the $2$-ramified $d_K=148$''. This
is the list of cubic fields with $d_K\le3000$ and $\nu\le3.7321$ irrespective of ramification;
besides $148$, the fields $316$, $404$, $564$ are also $2$-ramified and are therefore not covered
by the $2$-unramified threshold used there.
\item The limiting-variance theorem of \cite{PB} gives $\sqrt e$ as the lower bound for $\liminf\Var$;
Theorem~\ref{thm:lp} improves it.
\end{enumerate}

\section{Reproducibility}\label{sec:repro}
The repository contains:
\begin{itemize}
\item \texttt{src/gtest.gp}: the toolkit (field record, rigorous signs, the decision procedure
\texttt{sosdecide}, unit criterion, totally positive generators modulo $U_K^{2}$, the witness search
\texttt{gwitness2} with certification, and the invariants $\nu,\nu_2,t$);
\item \texttt{src/census.gp}, \texttt{src/runcensus.sh}: the census (one output line per field:
discriminant, polynomial, $h^+/h$, dyadic type, $\nu$, status, witness, norm, time, $\nu_2$, $h$,
regulator, normalised shape minimum, trace of the witness);
\item \texttt{src/run2o.gp}, \texttt{data/2o/}: $\Gcl$-witnesses for the fields left by \cite{TwoO};
\item \texttt{src/enum6.c} (copied from \cite{TwoO}), \texttt{src/varspec.sh},
\texttt{src/spectrum.gp}: the variance spectrum; \texttt{src/lp/varlp.py}, \texttt{src/lp/varlp\_cert.py}:
the linear programs and their interval-arithmetic verification;
\item \texttt{src/cyclo.gp}: real cyclotomic fields; \texttt{src/regime1.gp}: the ramified sextic
Case A test of \S\ref{sec:ramified}; \texttt{src/gen\_tables.py}: the tables of this note.
\end{itemize}
All runs used 4 cores; the full census takes about two hours, dominated by degree $8$ and by the
$193{,}467$ cubic fields.

\begin{thebibliography}{99}

\bibitem{CY} G.~Cherubini and P.~Yatsyna, \emph{Potential energy of totally positive algebraic
integers}, J. Math. Anal. Appl. \textbf{516} (2022), 126518.

\bibitem{Dubickas} A.~Dubickas, \emph{Algebraic integers with small potential energy}, J. Math.
Inequal. \textbf{17} (2023), 141--151.

\bibitem{DV} D.~S.~Dummit and J.~Voight, \emph{The $2$-Selmer group of a number field and
heuristics for narrow class groups and signature ranks of units}, Proc. Lond. Math. Soc. (3)
\textbf{117} (2018), 682--726.

\bibitem{KY} V.~Kala and P.~Yatsyna, \emph{Even better sums of squares over quintic and cyclotomic
fields}, arXiv:2402.03850.

\bibitem{LMFDB} The LMFDB Collaboration, \emph{The L-functions and modular forms database},
\url{https://www.lmfdb.org}.

\bibitem{MSRG} G.~Mantilla-Soler and C.~Rivera-Guaca, \emph{An introduction to Casimir pairings
and the shape of number fields}, arXiv:1907.09134.

\bibitem{OSS} B.~J.~Orloski, N.~T.~Sardari and A.~D.~Smith, \emph{New lower bounds for the
Schur--Siegel--Smyth trace problem}, Math. Comp. \textbf{94} (2025), 2005--2040.

\bibitem{PARI} The PARI~Group, \emph{PARI/GP, version 2.15.4}, Univ. Bordeaux, 2023,
\url{https://pari.math.u-bordeaux.fr/}.

\bibitem{PB} \emph{Even totally positive integers as sums of squares: the trace variance, a
signature law, and the real cyclotomic classification}, working notes (cited by the numbering of
the version supplied).

\bibitem{PD} \emph{A dyadic obstruction to sums of squares, and the real cyclotomic fields of
conductor $4m_0$}, working notes.

\bibitem{PR} \emph{Twisted capacity descent in the residual core of Kitaoka's conjecture, and a
relative trace--variance calculus}, working note.

\bibitem{Smith} A.~D.~Smith, \emph{Algebraic integers with conjugates in a prescribed
distribution}, Ann. of Math. (2) \textbf{200} (2024), 71--122.

\bibitem{TwoO} Companion code to \cite{KY}, repository \texttt{pav10/2O}, folder
\texttt{higher-degree} (notes \texttt{degrees6-7.tex}), \url{https://github.com/pav10/2O}.

\bibitem{Visser} R.~Visser, lists of totally real fields from the LMFDB, repository
\texttt{rvisser7/indecomposables}, \url{https://github.com/rvisser7/indecomposables}.

\end{thebibliography}
\end{document}
